\documentclass[11pt,a4paper]{article}
\usepackage[utf8]{inputenc}
\usepackage[T1]{fontenc}
\usepackage[catalan]{babel}
\usepackage{lmodern}
\usepackage[sfdefault,lining]{FiraSans}
\usepackage[a4paper,top=2.3cm,bottom=2.2cm,left=1.8cm,right=1.8cm,headheight=1.3cm,headsep=0.4cm,footskip=1cm]{geometry}
\usepackage{amsmath,amssymb}
\usepackage{icomma}
\usepackage{graphicx}
\usepackage[table]{xcolor}
\usepackage{tikz}
\usepackage{fontawesome5}
\usepackage{ccicons}
\usepackage{fancyhdr}
\usepackage{lastpage}
\usepackage{tasks}
\usepackage{needspace}
\usepackage[most]{tcolorbox}
\usepackage{array}
\usepackage[hidelinks,pdftitle={Fitxa d'exercicis de la unitat 9: Gràfics de funcions. 1r d'ESO},pdfsubject={Matemàtiques 1r ESO},pdfkeywords={coordenades cartesianes, eixos, quadrants, abscissa, ordenada, plànols, desplaçaments, simetries, taules de valors, gràfics, funcions, creixement, màxims i mínims, proporcionalitat directa, pendent, funció lineal, ordenada en l'origen}]{hyperref}

\definecolor{navy}{HTML}{1B2748}
\definecolor{accent}{HTML}{E8603A}
\definecolor{teal}{HTML}{1F9E90}
\definecolor{soft}{HTML}{EEF1F8}
\definecolor{muted}{HTML}{5D6685}

\setlength{\parindent}{0pt}
\setlength{\parskip}{3pt}
\settasks{label=\alph*), label-format=\bfseries\color{navy}, label-width=1.3em, item-indent=1.8em, column-sep=0.6em, after-item-skip=5pt, before-skip=3pt, after-skip=4pt}

% ---------- capçalera i peu ----------
\newcommand{\llicencia}{\href{https://creativecommons.org/licenses/by-nc-sa/4.0/deed.ca}{\ccbyncsa}\ \ Llicència Creative Commons Reconeixement-NoComercial-CompartirIgual 4.0 Internacional (CC BY-NC-SA 4.0)}
\pagestyle{fancy}
\fancyhf{}
\fancyhead[L]{\raisebox{-0.35\height}{\includegraphics[height=0.95cm]{logo-gva.png}}}
\fancyhead[C]{\small\color{muted}Fitxa d'exercicis · Unitat 9 · 1r d'ESO}
\fancyhead[R]{\raisebox{-0.35\height}{\includegraphics[height=0.95cm]{logo-vermellar.png}}}
\fancyfoot[L]{\scriptsize\color{muted}\llicencia}
\fancyfoot[R]{\small\color{muted}\thepage}
\renewcommand{\headrulewidth}{0.4pt}
\renewcommand{\headrule}{\color{muted!40}\hrule height \headrulewidth}
\fancypagestyle{primera}{\fancyhf{}\renewcommand{\headrulewidth}{0pt}
  \fancyfoot[L]{\scriptsize\color{muted}\llicencia}
  \fancyfoot[R]{\small\color{muted}\thepage}}

% ---------- nivells i exercicis ----------
\newcommand{\estrelles}[1]{{\color{accent}\ifcase#1\or\faStar\faStar[regular]\faStar[regular]\or\faStar\faStar\faStar[regular]\or\faStar\faStar\faStar\fi}}
\newcounter{exer}
\newcommand{\ex}[1]{\par\needspace{4\baselineskip}\medskip\stepcounter{exer}\noindent{\color{navy}\bfseries\theexer.}\ \ {\small\estrelles{#1}}\ \ }

\newcommand{\seccio}[2]{\par\needspace{10\baselineskip}\bigskip
  \begin{tcolorbox}[colback=navy,colframe=navy,coltext=white,arc=2mm,boxrule=0pt,left=3mm,right=3mm,top=1.2mm,bottom=1.2mm]
  {\large\bfseries #1\quad #2}
  \end{tcolorbox}\vspace{-2pt}}

\newcolumntype{C}[1]{>{\centering\arraybackslash}p{#1}}
\definecolor{violet}{HTML}{7C6CF0}

% quadre de recordatori (només a la fitxa d'exercicis)
\newcommand{\recorda}[1]{\par\smallskip\begin{tcolorbox}[colback=accent!7,colframe=accent!55,arc=2mm,boxrule=0.5pt,left=2.5mm,right=2.5mm,top=1mm,bottom=1mm]\small{\color{accent}\faLightbulb[regular]}\ \textbf{Recorda:}\ #1\end{tcolorbox}}
\newcommand{\res}[1]{#1}\newcommand{\rest}[1]{#1}

\def\fmneed{14}
% =====================================================================
% Capçalera FentMates de les fitxes: un color per a cada curs.
% Els colors són els de la pàgina inicial (LEVELS a assets/js/materials.js);
% el logo ve de _eines/logo/pdf/ (python3 _eines/logo/genera.py).
%
% Ús, al final del preàmbul de la unitat:
%   % =====================================================================
% Capçalera FentMates de les fitxes: un color per a cada curs.
% Els colors són els de la pàgina inicial (LEVELS a assets/js/materials.js);
% el logo ve de _eines/logo/pdf/ (python3 _eines/logo/genera.py).
%
% Ús, al final del preàmbul de la unitat:
%   % =====================================================================
% Capçalera FentMates de les fitxes: un color per a cada curs.
% Els colors són els de la pàgina inicial (LEVELS a assets/js/materials.js);
% el logo ve de _eines/logo/pdf/ (python3 _eines/logo/genera.py).
%
% Ús, al final del preàmbul de la unitat:
%   \input{../comu/fentmates.tex}
%   \fmcurs{1eso}
% i, al començament del document:
%   \capcalera{Matemàtiques}{1r ESO}{Fitxa d'exercicis · Unitat 2}{Divisibilitat}
% Opcional abans de \input: \def\fmneed{14} (espai mínim abans de cada secció).
% =====================================================================
\makeatletter
\def\fm@logodir{../../logo/pdf}
\providecommand{\fmneed}{10}
\def\fm@cursnom{}

% clau del curs, color i color del text damunt del color
\newcommand{\fm@curs}[3]{\@namedef{fm@color@#1}{#2}\@namedef{fm@text@#1}{#3}}
\fm@curs{1eso}{E6B84A}{navy}
\fm@curs{2eso}{3B82F6}{white}
\fm@curs{3eso}{FF7A50}{white}
\fm@curs{4eso-a}{A78BFA}{white}
\fm@curs{4eso-b}{6D28D9}{white}
\fm@curs{1batx-mi}{9B1C1C}{white}
\fm@curs{1batx-ccss}{EF6B6B}{white}
\fm@curs{2batx-mii}{1E6B47}{white}
\fm@curs{2batx-ccss}{2BB3A3}{white}

\newcommand{\fmcurs}[1]{%
  \@ifundefined{fm@color@#1}{\PackageError{fentmates}{Curs desconegut: #1}{}}{}%
  \edef\fm@tmp{\noexpand\definecolor{curs}{HTML}{\@nameuse{fm@color@#1}}%
               \noexpand\colorlet{curstext}{\@nameuse{fm@text@#1}}}\fm@tmp
  \colorlet{cursfosc}{curs!72!black}%
  \def\fm@logo{\fm@logodir/vertical-#1.pdf}%
  \def\fm@marca{\fm@logodir/marca-#1.pdf}}

% ---------- capçalera de la primera pàgina ----------
% #1 matèria (petit), #2 curs, #3 tipus de document i unitat, #4 títol
\newcommand{\capcalera}[4]{%
  \gdef\fm@cursnom{#2}%
  \par\noindent
  \begin{tcolorbox}[colback=curs,colframe=curs,boxrule=0pt,arc=2.2mm,width=3.5cm,height=2.15cm,
      valign=center,left=3mm,right=2mm,top=1mm,bottom=1mm,nobeforeafter,box align=center]
    \color{curstext}{\footnotesize #1}\par\vspace{1pt}{\fontsize{19}{21}\selectfont\bfseries #2\par}%
  \end{tcolorbox}\hspace{5mm}%
  \begin{minipage}[c]{\dimexpr\linewidth-3.5cm-5mm-2.9cm\relax}
    {\color{cursfosc}\small\bfseries\MakeUppercase{#3}}\par\vspace{2pt}
    {\color{navy}\fontsize{20}{23}\selectfont\bfseries #4\par}
  \end{minipage}\hfill
  \begin{minipage}[c]{2.6cm}\raggedleft\includegraphics[width=2.6cm]{\fm@logo}\end{minipage}\par
  \vspace{2pt}}

% ---------- barres de secció del color del curs ----------
\renewcommand{\seccio}[2]{\par\needspace{\fmneed\baselineskip}\bigskip
  \begin{tcolorbox}[colback=curs,colframe=curs,coltext=curstext,arc=2mm,boxrule=0pt,left=3mm,right=3mm,top=1.2mm,bottom=1.2mm]
  {\large\bfseries #1\quad #2}
  \end{tcolorbox}\vspace{-2pt}}

% ---------- estrelles de dificultat del color del curs (a Batxillerat no se'n posen) ----------
\@ifundefined{estrelles}{}{%
  \renewcommand{\estrelles}[1]{{\color{curs}\ifcase#1\or\faStar\faStar[regular]\faStar[regular]\or\faStar\faStar\faStar[regular]\or\faStar\faStar\faStar\fi}}}

% ---------- capçalera i peu de les pàgines ----------
% primera pàgina: logos institucionals al peu
\fancypagestyle{primera}{\fancyhf{}\renewcommand{\headrulewidth}{0pt}
  \fancyfoot[L]{\raisebox{-0.32\height}{\includegraphics[height=0.78cm]{logo-gva.png}}\hspace{5mm}%
    \raisebox{-0.32\height}{\includegraphics[height=0.78cm]{logo-vermellar.png}}}
  \fancyfoot[R]{\scriptsize\color{muted}\href{https://creativecommons.org/licenses/by-nc-sa/4.0/deed.ca}{\ccbyncsa}\ \ CC BY-NC-SA 4.0 · fentmates.org\qquad\small\thepage}}
% la resta: el símbol a l'esquerra i el curs a la dreta, amb el títol al centre com abans
\fancyhead[L]{\raisebox{-0.3\height}{\includegraphics[height=0.62cm]{\fm@marca}}}
\fancyhead[R]{\small\bfseries\color{cursfosc}\fm@cursnom}
\renewcommand{\headrule}{\color{curs!45}\hrule height \headrulewidth}
\makeatother

%   \fmcurs{1eso}
% i, al començament del document:
%   \capcalera{Matemàtiques}{1r ESO}{Fitxa d'exercicis · Unitat 2}{Divisibilitat}
% Opcional abans de \input: \def\fmneed{14} (espai mínim abans de cada secció).
% =====================================================================
\makeatletter
\def\fm@logodir{../../logo/pdf}
\providecommand{\fmneed}{10}
\def\fm@cursnom{}

% clau del curs, color i color del text damunt del color
\newcommand{\fm@curs}[3]{\@namedef{fm@color@#1}{#2}\@namedef{fm@text@#1}{#3}}
\fm@curs{1eso}{E6B84A}{navy}
\fm@curs{2eso}{3B82F6}{white}
\fm@curs{3eso}{FF7A50}{white}
\fm@curs{4eso-a}{A78BFA}{white}
\fm@curs{4eso-b}{6D28D9}{white}
\fm@curs{1batx-mi}{9B1C1C}{white}
\fm@curs{1batx-ccss}{EF6B6B}{white}
\fm@curs{2batx-mii}{1E6B47}{white}
\fm@curs{2batx-ccss}{2BB3A3}{white}

\newcommand{\fmcurs}[1]{%
  \@ifundefined{fm@color@#1}{\PackageError{fentmates}{Curs desconegut: #1}{}}{}%
  \edef\fm@tmp{\noexpand\definecolor{curs}{HTML}{\@nameuse{fm@color@#1}}%
               \noexpand\colorlet{curstext}{\@nameuse{fm@text@#1}}}\fm@tmp
  \colorlet{cursfosc}{curs!72!black}%
  \def\fm@logo{\fm@logodir/vertical-#1.pdf}%
  \def\fm@marca{\fm@logodir/marca-#1.pdf}}

% ---------- capçalera de la primera pàgina ----------
% #1 matèria (petit), #2 curs, #3 tipus de document i unitat, #4 títol
\newcommand{\capcalera}[4]{%
  \gdef\fm@cursnom{#2}%
  \par\noindent
  \begin{tcolorbox}[colback=curs,colframe=curs,boxrule=0pt,arc=2.2mm,width=3.5cm,height=2.15cm,
      valign=center,left=3mm,right=2mm,top=1mm,bottom=1mm,nobeforeafter,box align=center]
    \color{curstext}{\footnotesize #1}\par\vspace{1pt}{\fontsize{19}{21}\selectfont\bfseries #2\par}%
  \end{tcolorbox}\hspace{5mm}%
  \begin{minipage}[c]{\dimexpr\linewidth-3.5cm-5mm-2.9cm\relax}
    {\color{cursfosc}\small\bfseries\MakeUppercase{#3}}\par\vspace{2pt}
    {\color{navy}\fontsize{20}{23}\selectfont\bfseries #4\par}
  \end{minipage}\hfill
  \begin{minipage}[c]{2.6cm}\raggedleft\includegraphics[width=2.6cm]{\fm@logo}\end{minipage}\par
  \vspace{2pt}}

% ---------- barres de secció del color del curs ----------
\renewcommand{\seccio}[2]{\par\needspace{\fmneed\baselineskip}\bigskip
  \begin{tcolorbox}[colback=curs,colframe=curs,coltext=curstext,arc=2mm,boxrule=0pt,left=3mm,right=3mm,top=1.2mm,bottom=1.2mm]
  {\large\bfseries #1\quad #2}
  \end{tcolorbox}\vspace{-2pt}}

% ---------- estrelles de dificultat del color del curs (a Batxillerat no se'n posen) ----------
\@ifundefined{estrelles}{}{%
  \renewcommand{\estrelles}[1]{{\color{curs}\ifcase#1\or\faStar\faStar[regular]\faStar[regular]\or\faStar\faStar\faStar[regular]\or\faStar\faStar\faStar\fi}}}

% ---------- capçalera i peu de les pàgines ----------
% primera pàgina: logos institucionals al peu
\fancypagestyle{primera}{\fancyhf{}\renewcommand{\headrulewidth}{0pt}
  \fancyfoot[L]{\raisebox{-0.32\height}{\includegraphics[height=0.78cm]{logo-gva.png}}\hspace{5mm}%
    \raisebox{-0.32\height}{\includegraphics[height=0.78cm]{logo-vermellar.png}}}
  \fancyfoot[R]{\scriptsize\color{muted}\href{https://creativecommons.org/licenses/by-nc-sa/4.0/deed.ca}{\ccbyncsa}\ \ CC BY-NC-SA 4.0 · fentmates.org\qquad\small\thepage}}
% la resta: el símbol a l'esquerra i el curs a la dreta, amb el títol al centre com abans
\fancyhead[L]{\raisebox{-0.3\height}{\includegraphics[height=0.62cm]{\fm@marca}}}
\fancyhead[R]{\small\bfseries\color{cursfosc}\fm@cursnom}
\renewcommand{\headrule}{\color{curs!45}\hrule height \headrulewidth}
\makeatother

%   \fmcurs{1eso}
% i, al començament del document:
%   \capcalera{Matemàtiques}{1r ESO}{Fitxa d'exercicis · Unitat 2}{Divisibilitat}
% Opcional abans de \input: \def\fmneed{14} (espai mínim abans de cada secció).
% =====================================================================
\makeatletter
\def\fm@logodir{../../logo/pdf}
\providecommand{\fmneed}{10}
\def\fm@cursnom{}

% clau del curs, color i color del text damunt del color
\newcommand{\fm@curs}[3]{\@namedef{fm@color@#1}{#2}\@namedef{fm@text@#1}{#3}}
\fm@curs{1eso}{E6B84A}{navy}
\fm@curs{2eso}{3B82F6}{white}
\fm@curs{3eso}{FF7A50}{white}
\fm@curs{4eso-a}{A78BFA}{white}
\fm@curs{4eso-b}{6D28D9}{white}
\fm@curs{1batx-mi}{9B1C1C}{white}
\fm@curs{1batx-ccss}{EF6B6B}{white}
\fm@curs{2batx-mii}{1E6B47}{white}
\fm@curs{2batx-ccss}{2BB3A3}{white}

\newcommand{\fmcurs}[1]{%
  \@ifundefined{fm@color@#1}{\PackageError{fentmates}{Curs desconegut: #1}{}}{}%
  \edef\fm@tmp{\noexpand\definecolor{curs}{HTML}{\@nameuse{fm@color@#1}}%
               \noexpand\colorlet{curstext}{\@nameuse{fm@text@#1}}}\fm@tmp
  \colorlet{cursfosc}{curs!72!black}%
  \def\fm@logo{\fm@logodir/vertical-#1.pdf}%
  \def\fm@marca{\fm@logodir/marca-#1.pdf}}

% ---------- capçalera de la primera pàgina ----------
% #1 matèria (petit), #2 curs, #3 tipus de document i unitat, #4 títol
\newcommand{\capcalera}[4]{%
  \gdef\fm@cursnom{#2}%
  \par\noindent
  \begin{tcolorbox}[colback=curs,colframe=curs,boxrule=0pt,arc=2.2mm,width=3.5cm,height=2.15cm,
      valign=center,left=3mm,right=2mm,top=1mm,bottom=1mm,nobeforeafter,box align=center]
    \color{curstext}{\footnotesize #1}\par\vspace{1pt}{\fontsize{19}{21}\selectfont\bfseries #2\par}%
  \end{tcolorbox}\hspace{5mm}%
  \begin{minipage}[c]{\dimexpr\linewidth-3.5cm-5mm-2.9cm\relax}
    {\color{cursfosc}\small\bfseries\MakeUppercase{#3}}\par\vspace{2pt}
    {\color{navy}\fontsize{20}{23}\selectfont\bfseries #4\par}
  \end{minipage}\hfill
  \begin{minipage}[c]{2.6cm}\raggedleft\includegraphics[width=2.6cm]{\fm@logo}\end{minipage}\par
  \vspace{2pt}}

% ---------- barres de secció del color del curs ----------
\renewcommand{\seccio}[2]{\par\needspace{\fmneed\baselineskip}\bigskip
  \begin{tcolorbox}[colback=curs,colframe=curs,coltext=curstext,arc=2mm,boxrule=0pt,left=3mm,right=3mm,top=1.2mm,bottom=1.2mm]
  {\large\bfseries #1\quad #2}
  \end{tcolorbox}\vspace{-2pt}}

% ---------- estrelles de dificultat del color del curs (a Batxillerat no se'n posen) ----------
\@ifundefined{estrelles}{}{%
  \renewcommand{\estrelles}[1]{{\color{curs}\ifcase#1\or\faStar\faStar[regular]\faStar[regular]\or\faStar\faStar\faStar[regular]\or\faStar\faStar\faStar\fi}}}

% ---------- capçalera i peu de les pàgines ----------
% primera pàgina: logos institucionals al peu
\fancypagestyle{primera}{\fancyhf{}\renewcommand{\headrulewidth}{0pt}
  \fancyfoot[L]{\raisebox{-0.32\height}{\includegraphics[height=0.78cm]{logo-gva.png}}\hspace{5mm}%
    \raisebox{-0.32\height}{\includegraphics[height=0.78cm]{logo-vermellar.png}}}
  \fancyfoot[R]{\scriptsize\color{muted}\href{https://creativecommons.org/licenses/by-nc-sa/4.0/deed.ca}{\ccbyncsa}\ \ CC BY-NC-SA 4.0 · fentmates.org\qquad\small\thepage}}
% la resta: el símbol a l'esquerra i el curs a la dreta, amb el títol al centre com abans
\fancyhead[L]{\raisebox{-0.3\height}{\includegraphics[height=0.62cm]{\fm@marca}}}
\fancyhead[R]{\small\bfseries\color{cursfosc}\fm@cursnom}
\renewcommand{\headrule}{\color{curs!45}\hrule height \headrulewidth}
\makeatother

\fmcurs{1eso}

\begin{document}
\thispagestyle{primera}

\capcalera{Matemàtiques}{1r ESO}{Fitxa d'exercicis · Unitat 9}{Gràfics de funcions}
\vspace{10pt}
Nom i cognoms: \rule{8.4cm}{0.4pt}\hfill Grup: \rule{1.6cm}{0.4pt}\hfill Data: \rule{2.2cm}{0.4pt}
\vspace{8pt}

\begin{tcolorbox}[colback=soft,colframe=soft,arc=2mm,left=3mm,right=3mm,top=1.5mm,bottom=1.5mm]
\small Exercicis graduats per nivells de dificultat:\quad \estrelles{1}\ bàsic \qquad \estrelles{2}\ mitjà \qquad \estrelles{3}\ avançat.\\[2pt] Resol-los al quadern mostrant tots els passos. Per a dibuixar punts, rectes i gràfics fes servir paper quadriculat i escriu sempre el nom i la unitat de cada eix.
\end{tcolorbox}

% =====================================================
\seccio{1}{El pla cartesià. Coordenades d'un punt}

\ex{1} Escriu les coordenades dels punts del dibuix:\par\smallskip\begin{center}\begin{tikzpicture}[x=0.42cm,y=0.42cm,baseline=(current bounding box.center),line cap=round,line join=round]
\draw[muted!30,line width=0.3pt] (-5,-5)--(-5,5);
\draw[muted!30,line width=0.3pt] (-4,-5)--(-4,5);
\draw[muted!30,line width=0.3pt] (-3,-5)--(-3,5);
\draw[muted!30,line width=0.3pt] (-2,-5)--(-2,5);
\draw[muted!30,line width=0.3pt] (-1,-5)--(-1,5);
\draw[muted!30,line width=0.3pt] (0,-5)--(0,5);
\draw[muted!30,line width=0.3pt] (1,-5)--(1,5);
\draw[muted!30,line width=0.3pt] (2,-5)--(2,5);
\draw[muted!30,line width=0.3pt] (3,-5)--(3,5);
\draw[muted!30,line width=0.3pt] (4,-5)--(4,5);
\draw[muted!30,line width=0.3pt] (5,-5)--(5,5);
\draw[muted!30,line width=0.3pt] (-5,-5)--(5,-5);
\draw[muted!30,line width=0.3pt] (-5,-4)--(5,-4);
\draw[muted!30,line width=0.3pt] (-5,-3)--(5,-3);
\draw[muted!30,line width=0.3pt] (-5,-2)--(5,-2);
\draw[muted!30,line width=0.3pt] (-5,-1)--(5,-1);
\draw[muted!30,line width=0.3pt] (-5,0)--(5,0);
\draw[muted!30,line width=0.3pt] (-5,1)--(5,1);
\draw[muted!30,line width=0.3pt] (-5,2)--(5,2);
\draw[muted!30,line width=0.3pt] (-5,3)--(5,3);
\draw[muted!30,line width=0.3pt] (-5,4)--(5,4);
\draw[muted!30,line width=0.3pt] (-5,5)--(5,5);
\draw[-latex,muted,line width=0.6pt] (-5,0)--(5.3,0) node[right,font=\scriptsize,text=navy]{x};
\draw[-latex,muted,line width=0.6pt] (0,-5)--(0,5.3) node[above,font=\scriptsize,text=navy]{y};
\node[font=\tiny,text=muted,below,inner sep=1.5pt] at (-4,0) {$-4$};
\node[font=\tiny,text=muted,below,inner sep=1.5pt] at (-3,0) {$-3$};
\node[font=\tiny,text=muted,below,inner sep=1.5pt] at (-2,0) {$-2$};
\node[font=\tiny,text=muted,below,inner sep=1.5pt] at (-1,0) {$-1$};
\node[font=\tiny,text=muted,below,inner sep=1.5pt] at (1,0) {$1$};
\node[font=\tiny,text=muted,below,inner sep=1.5pt] at (2,0) {$2$};
\node[font=\tiny,text=muted,below,inner sep=1.5pt] at (3,0) {$3$};
\node[font=\tiny,text=muted,below,inner sep=1.5pt] at (4,0) {$4$};
\node[font=\tiny,text=muted,left,inner sep=1.5pt] at (0,-4) {$-4$};
\node[font=\tiny,text=muted,left,inner sep=1.5pt] at (0,-3) {$-3$};
\node[font=\tiny,text=muted,left,inner sep=1.5pt] at (0,-2) {$-2$};
\node[font=\tiny,text=muted,left,inner sep=1.5pt] at (0,-1) {$-1$};
\node[font=\tiny,text=muted,left,inner sep=1.5pt] at (0,1) {$1$};
\node[font=\tiny,text=muted,left,inner sep=1.5pt] at (0,2) {$2$};
\node[font=\tiny,text=muted,left,inner sep=1.5pt] at (0,3) {$3$};
\node[font=\tiny,text=muted,left,inner sep=1.5pt] at (0,4) {$4$};
\node[font=\tiny,text=muted,below left,inner sep=1.5pt] at (0,0) {$0$};
\fill[navy] (2,3) circle (2.2pt) node[above right,font=\scriptsize,inner sep=1.5pt]{$A$};
\fill[navy] (-4,1) circle (2.2pt) node[above right,font=\scriptsize,inner sep=1.5pt]{$B$};
\fill[navy] (-3,-2) circle (2.2pt) node[above right,font=\scriptsize,inner sep=1.5pt]{$C$};
\fill[navy] (3,-4) circle (2.2pt) node[above right,font=\scriptsize,inner sep=1.5pt]{$D$};
\fill[navy] (0,2) circle (2.2pt) node[above right,font=\scriptsize,inner sep=1.5pt]{$E$};
\fill[navy] (4,0) circle (2.2pt) node[above right,font=\scriptsize,inner sep=1.5pt]{$F$};
\fill[navy] (-1,4) circle (2.2pt) node[above right,font=\scriptsize,inner sep=1.5pt]{$G$};
\fill[navy] (-2,0) circle (2.2pt) node[above right,font=\scriptsize,inner sep=1.5pt]{$H$};
\end{tikzpicture}\end{center}\recorda{Les coordenades d'un punt s'escriuen $P(x,\ y)$: primer l'\textbf{abscissa} $x$ (horitzontal: dreta $+$, esquerra $-$) i després l'\textbf{ordenada} $y$ (vertical: amunt $+$, avall $-$).}

\ex{1} Dibuixa en uns eixos de coordenades els punts $P(3,\ 1)$, $Q(-2,\ 4)$, $R(-4,\ -3)$, $S(1,\ -2)$, $T(0,\ -4)$ i $U(5,\ 0)$.

\ex{1} Sense dibuixar-los, digues en quin quadrant o sobre quin eix està cada punt:\begin{tasks}(4)
  \task $(6,\ -2)$
  \task $(-3,\ -8)$
  \task $(-1,\ 7)$
  \task $(0,\ -5)$
  \task $(9,\ 0)$
  \task $(4,\ 4)$
  \task $(-6,\ 0)$
  \task $(-2,\ -2)$
\end{tasks}\recorda{Signes: quadrant I $(+,\,+)$, II $(-,\,+)$, III $(-,\,-)$, IV $(+,\,-)$. Si $y=0$, el punt és de l'eix $X$; si $x=0$, de l'eix $Y$.}

\ex{1} Escriu les coordenades del punt:\begin{tasks}(1)
  \task que està 3 unitats a la dreta de l'eix $Y$ i 5 per damunt de l'eix $X$;
  \task que està 2 unitats a l'esquerra de l'eix $Y$ i 4 per davall de l'eix $X$;
  \task que està sobre l'eix $X$, 6 unitats a l'esquerra de l'origen;
  \task que està sobre l'eix $Y$, 1 unitat per damunt de l'origen.
\end{tasks}

\ex{2} Dibuixa els punts $(-2,\ 4)$ i $(4,\ -2)$. Són el mateix punt? En quin quadrant està cadascun?

\ex{2} Respon:\begin{tasks}(1)
  \task Quin punt del quadrant II està a 4 unitats de cada eix?
  \task Quins punts de l'eix $X$ estan a 3 unitats de l'origen?
  \task Quin punt del quadrant III té l'abscissa $-5$ i l'ordenada la meitat de l'abscissa?
\end{tasks}

\ex{2} El punt $(a,\ b)$ està en el quadrant IV. En quin quadrant estan els punts $(-a,\ b)$, \ $(a,\ -b)$ \ i \ $(b,\ a)$?

\ex{3} Dibuixa tots els punts de coordenades enteres entre $-4$ i $4$ que tenen l'ordenada igual a l'abscissa. Fes el mateix amb els que tenen l'ordenada igual a 2. Què observes?

\ex{3} Dibuixa els punts $A(-3,\ -2)$, $B(1,\ -2)$ i $C(1,\ 2)$. Troba el punt $D$ perquè $ABCD$ siga un quadrat i calcula les coordenades del centre del quadrat.


% =====================================================
\seccio{2}{Aplicacions de les coordenades}

\ex{1} En aquest plànol d'un poble, la bandera és la plaça (l'origen) i cada quadre és una illa de cases. Escriu les coordenades de l'escola ($E$), el parc ($P$), l'hospital ($H$) i el mercat ($M$). Què hi ha en $(0,\ -1)$? I en $(-4,\ 0)$?\par\smallskip\begin{center}\begin{tikzpicture}[x=0.42cm,y=0.42cm,baseline=(current bounding box.center),line cap=round,line join=round]
\draw[muted!30,line width=0.3pt] (-5,-5)--(-5,5);
\draw[muted!30,line width=0.3pt] (-4,-5)--(-4,5);
\draw[muted!30,line width=0.3pt] (-3,-5)--(-3,5);
\draw[muted!30,line width=0.3pt] (-2,-5)--(-2,5);
\draw[muted!30,line width=0.3pt] (-1,-5)--(-1,5);
\draw[muted!30,line width=0.3pt] (0,-5)--(0,5);
\draw[muted!30,line width=0.3pt] (1,-5)--(1,5);
\draw[muted!30,line width=0.3pt] (2,-5)--(2,5);
\draw[muted!30,line width=0.3pt] (3,-5)--(3,5);
\draw[muted!30,line width=0.3pt] (4,-5)--(4,5);
\draw[muted!30,line width=0.3pt] (5,-5)--(5,5);
\draw[muted!30,line width=0.3pt] (-5,-5)--(5,-5);
\draw[muted!30,line width=0.3pt] (-5,-4)--(5,-4);
\draw[muted!30,line width=0.3pt] (-5,-3)--(5,-3);
\draw[muted!30,line width=0.3pt] (-5,-2)--(5,-2);
\draw[muted!30,line width=0.3pt] (-5,-1)--(5,-1);
\draw[muted!30,line width=0.3pt] (-5,0)--(5,0);
\draw[muted!30,line width=0.3pt] (-5,1)--(5,1);
\draw[muted!30,line width=0.3pt] (-5,2)--(5,2);
\draw[muted!30,line width=0.3pt] (-5,3)--(5,3);
\draw[muted!30,line width=0.3pt] (-5,4)--(5,4);
\draw[muted!30,line width=0.3pt] (-5,5)--(5,5);
\draw[-latex,muted,line width=0.6pt] (-5,0)--(5.3,0) node[right,font=\scriptsize,text=navy]{x};
\draw[-latex,muted,line width=0.6pt] (0,-5)--(0,5.3) node[above,font=\scriptsize,text=navy]{y};
\node[font=\tiny,text=muted,below,inner sep=1.5pt] at (-4,0) {$-4$};
\node[font=\tiny,text=muted,below,inner sep=1.5pt] at (-3,0) {$-3$};
\node[font=\tiny,text=muted,below,inner sep=1.5pt] at (-2,0) {$-2$};
\node[font=\tiny,text=muted,below,inner sep=1.5pt] at (-1,0) {$-1$};
\node[font=\tiny,text=muted,below,inner sep=1.5pt] at (1,0) {$1$};
\node[font=\tiny,text=muted,below,inner sep=1.5pt] at (2,0) {$2$};
\node[font=\tiny,text=muted,below,inner sep=1.5pt] at (3,0) {$3$};
\node[font=\tiny,text=muted,below,inner sep=1.5pt] at (4,0) {$4$};
\node[font=\tiny,text=muted,left,inner sep=1.5pt] at (0,-4) {$-4$};
\node[font=\tiny,text=muted,left,inner sep=1.5pt] at (0,-3) {$-3$};
\node[font=\tiny,text=muted,left,inner sep=1.5pt] at (0,-2) {$-2$};
\node[font=\tiny,text=muted,left,inner sep=1.5pt] at (0,-1) {$-1$};
\node[font=\tiny,text=muted,left,inner sep=1.5pt] at (0,1) {$1$};
\node[font=\tiny,text=muted,left,inner sep=1.5pt] at (0,2) {$2$};
\node[font=\tiny,text=muted,left,inner sep=1.5pt] at (0,3) {$3$};
\node[font=\tiny,text=muted,left,inner sep=1.5pt] at (0,4) {$4$};
\node[font=\tiny,text=muted,below left,inner sep=1.5pt] at (0,0) {$0$};
\fill[accent] (2,3) circle (2.2pt) node[above right,font=\scriptsize,inner sep=1.5pt]{\textbf{E}};
\fill[accent] (-3,2) circle (2.2pt) node[above right,font=\scriptsize,inner sep=1.5pt]{\textbf{P}};
\fill[accent] (4,-2) circle (2.2pt) node[above right,font=\scriptsize,inner sep=1.5pt]{\textbf{H}};
\fill[accent] (-2,-3) circle (2.2pt) node[above right,font=\scriptsize,inner sep=1.5pt]{\textbf{M}};
\fill[accent] (0,-1) circle (2.2pt) node[above right,font=\scriptsize,inner sep=1.5pt]{\textbf{B}};
\fill[accent] (-4,0) circle (2.2pt) node[above right,font=\scriptsize,inner sep=1.5pt]{\textbf{C}};
\node[font=\scriptsize,above right,text=muted] at (0.25,0.25) {\faFlag};
\end{tikzpicture}\end{center}

\ex{1} Dibuixa el polígon de vèrtexs $A(1,\ 1)$, $B(6,\ 1)$, $C(6,\ 4)$ i $D(1,\ 4)$. Quin polígon és? Calcula'n el perímetre i l'àrea (la unitat és el cm).\recorda{Si dos punts tenen la mateixa $y$, la distància entre ells és la diferència de les $x$; si tenen la mateixa $x$, la diferència de les $y$.}

\ex{1} Calcula la distància entre els punts:\begin{tasks}(2)
  \task $(-2,\ 3)$ i $(5,\ 3)$
  \task $(4,\ -1)$ i $(4,\ 6)$
  \task $(-6,\ -2)$ i $(-1,\ -2)$
  \task $(0,\ -4)$ i $(0,\ 3)$
\end{tasks}

\ex{2} Un robot ix de l'origen i fa aquests moviments: 4 a la dreta, 3 amunt, 6 a l'esquerra, 4 avall i 3 a la dreta. On acaba? Dibuixa'n el camí.

\ex{2} Desplaça el triangle de vèrtexs $(-4,\ 1)$, $(-1,\ 1)$ i $(-2,\ 4)$ 5 unitats a la dreta i 4 cap avall. Escriu els vèrtexs nous i dibuixa els dos triangles.\recorda{Desplaçar $a$ a la dreta i $b$ amunt: $(x,\ y)\to(x+a,\ y+b)$ (si va a l'esquerra o avall, restem).}

\ex{2} Completa la taula amb els simètrics de cada punt respecte dels eixos:\par\smallskip\begin{center}\renewcommand{\arraystretch}{1.3}\begin{tabular}{|c|C{2.2cm}|C{2.2cm}|}\hline Punt & Simètric respecte de l'eix $X$ & Simètric respecte de l'eix $Y$\\\hline $(3,\ 5)$ & & \\\hline $(-2,\ 4)$ & & \\\hline $(-1,\ -6)$ & & \\\hline $(5,\ -3)$ & & \\\hline\end{tabular}\end{center}\recorda{Respecte de l'eix $X$ canvia el signe de la $y$: $(x,\ -y)$. Respecte de l'eix $Y$ canvia el signe de la $x$: $(-x,\ y)$.}

\ex{2} Tres vèrtexs d'un rectangle són $A(-2,\ 1)$, $B(3,\ 1)$ i $C(3,\ -2)$. Troba el quart vèrtex $D$ i calcula'n l'àrea i el perímetre.

\ex{3} Dibuixa el triangle de vèrtexs $A(-2,\ -1)$, $B(4,\ -1)$ i $C(1,\ 3)$ i calcula'n l'àrea.

\ex{3} Dibuixa el pentàgon de vèrtexs $(0,\ 0)$, $(4,\ 0)$, $(4,\ 2)$, $(2,\ 4)$ i $(0,\ 2)$ i calcula'n l'àrea.


% =====================================================
\seccio{3}{Relacions entre magnituds: taules i gràfics}

\ex{1} Un panet val 0,40\,€. Completa la taula, dibuixa el gràfic i digues si s'han d'unir els punts.\par\smallskip\begin{center}$\begin{array}{|c|C{0.75cm}C{0.75cm}C{0.75cm}C{0.75cm}C{0.75cm}C{0.75cm}|}\hline \text{Panets}&1&2&3&4&5&6\\\hline \text{Preu (€)}&\phantom{00}&\phantom{00}&\phantom{00}&\phantom{00}&\phantom{00}&\phantom{00}\\\hline\end{array}$\end{center}\recorda{Cada parell $(x,\ y)$ de la taula és un punt del gràfic. Unim els punts només si també tenen sentit els valors intermedis.}

\ex{1} Digues quina és la variable independent i quina la dependent:\begin{tasks}(1)
  \task el temps que està oberta una aixeta i els litres d'aigua que ixen;
  \task el nombre de fotocòpies i el preu que paguem;
  \task l'hora del dia i la temperatura;
  \task els km que fa un cotxe i la gasolina que li queda.
\end{tasks}

\ex{1} Completa la taula de la funció $y=2x+3$, escriu els punts i dibuixa'ls.\par\smallskip\begin{center}$\begin{array}{|c|C{0.75cm}C{0.75cm}C{0.75cm}C{0.75cm}C{0.75cm}|}\hline x&0&1&2&3&4\\\hline y&\phantom{00}&\phantom{00}&\phantom{00}&\phantom{00}&\phantom{00}\\\hline\end{array}$\end{center}

\ex{1} La taula dona la temperatura d'un dia de primavera. Dibuixa el gràfic (unint els punts) i digues a quina hora va fer més calor i a quina més fred.\par\smallskip\begin{center}$\begin{array}{|c|ccccccccc|}\hline \text{Hora}&0&3&6&9&12&15&18&21&24\\\hline ^\circ\text{C}&8&6&7&12&16&18&15&11&9\\\hline\end{array}$\end{center}

\ex{2} Quines d'aquestes taules corresponen a una funció? Explica per què.\par\smallskip\begin{center}$\begin{array}{|c|cccc|}\hline x&1&2&3&4\\\hline y&5&10&15&20\\\hline\end{array}$\quad $\begin{array}{|c|cccc|}\hline x&1&1&2&3\\\hline y&4&6&8&9\\\hline\end{array}$\quad $\begin{array}{|c|cccc|}\hline x&0&2&4&6\\\hline y&3&3&3&3\\\hline\end{array}$\end{center}

\ex{2} Troba la fórmula de cada taula:\par\smallskip\begin{center}$\begin{array}{|c|cccc|}\hline x&0&1&2&3\\\hline y&5&7&9&11\\\hline\end{array}$\quad $\begin{array}{|c|cccc|}\hline x&1&2&3&4\\\hline y&10&20&30&40\\\hline\end{array}$\quad $\begin{array}{|c|cccc|}\hline x&0&1&2&3\\\hline y&9&8&7&6\\\hline\end{array}$\end{center}

\ex{2} Una tarifa de mòbil costa 10\,€ al mes més 3\,€ per cada GB de dades. Escriu la fórmula del preu $y$ si gastem $x$ GB i fes una taula per a 0, 1, 2, 5 i 10 GB.

\ex{3} Fes una taula del perímetre $P$ i una de l'àrea $A$ d'un quadrat de costat $l=1,\ 2,\ 3,\ 4,\ 5$\,cm. Escriu-ne les fórmules i dibuixa els gràfics. Quin és una recta?

\ex{3} Volem fer rectangles de 20\,cm de perímetre amb costats de cm enters. Fes una taula amb la base $x$ i l'altura $y$, escriu-ne la fórmula i dibuixa el gràfic. És creixent o decreixent?


% =====================================================
\seccio{4}{Interpretació de gràfics}

\ex{1} El gràfic mostra l'excursió en bici de la Núria (distància a casa al llarg del temps).\par\smallskip\begin{center}\begin{tikzpicture}[x=0.75cm,y=0.3cm,baseline=(current bounding box.center),line cap=round,line join=round]
\draw[muted!30,line width=0.3pt] (0,0)--(0,16);
\draw[muted!30,line width=0.3pt] (1,0)--(1,16);
\draw[muted!30,line width=0.3pt] (2,0)--(2,16);
\draw[muted!30,line width=0.3pt] (3,0)--(3,16);
\draw[muted!30,line width=0.3pt] (4,0)--(4,16);
\draw[muted!30,line width=0.3pt] (5,0)--(5,16);
\draw[muted!30,line width=0.3pt] (6,0)--(6,16);
\draw[muted!30,line width=0.3pt] (7,0)--(7,16);
\draw[muted!30,line width=0.3pt] (0,0)--(7.5,0);
\draw[muted!30,line width=0.3pt] (0,2)--(7.5,2);
\draw[muted!30,line width=0.3pt] (0,4)--(7.5,4);
\draw[muted!30,line width=0.3pt] (0,6)--(7.5,6);
\draw[muted!30,line width=0.3pt] (0,8)--(7.5,8);
\draw[muted!30,line width=0.3pt] (0,10)--(7.5,10);
\draw[muted!30,line width=0.3pt] (0,12)--(7.5,12);
\draw[muted!30,line width=0.3pt] (0,14)--(7.5,14);
\draw[muted!30,line width=0.3pt] (0,16)--(7.5,16);
\draw[-latex,muted,line width=0.6pt] (0,0)--(7.8,0) node[right,font=\scriptsize,text=navy]{temps (h)};
\draw[-latex,muted,line width=0.6pt] (0,0)--(0,16.6) node[above,font=\scriptsize,text=navy]{distància (km)};
\node[font=\tiny,text=muted,below,inner sep=1.5pt] at (1,0) {$1$};
\node[font=\tiny,text=muted,below,inner sep=1.5pt] at (2,0) {$2$};
\node[font=\tiny,text=muted,below,inner sep=1.5pt] at (3,0) {$3$};
\node[font=\tiny,text=muted,below,inner sep=1.5pt] at (4,0) {$4$};
\node[font=\tiny,text=muted,below,inner sep=1.5pt] at (5,0) {$5$};
\node[font=\tiny,text=muted,below,inner sep=1.5pt] at (6,0) {$6$};
\node[font=\tiny,text=muted,below,inner sep=1.5pt] at (7,0) {$7$};
\node[font=\tiny,text=muted,left,inner sep=1.5pt] at (0,2) {$2$};
\node[font=\tiny,text=muted,left,inner sep=1.5pt] at (0,4) {$4$};
\node[font=\tiny,text=muted,left,inner sep=1.5pt] at (0,6) {$6$};
\node[font=\tiny,text=muted,left,inner sep=1.5pt] at (0,8) {$8$};
\node[font=\tiny,text=muted,left,inner sep=1.5pt] at (0,10) {$10$};
\node[font=\tiny,text=muted,left,inner sep=1.5pt] at (0,12) {$12$};
\node[font=\tiny,text=muted,left,inner sep=1.5pt] at (0,14) {$14$};
\node[font=\tiny,text=muted,below left,inner sep=1.5pt] at (0,0) {$0$};
\draw[accent,line width=1.1pt] (0,0)--(1,6)--(2,6)--(4,14)--(5,14)--(7,0);
\fill[navy] (0,0) circle (2.2pt);
\fill[navy] (1,6) circle (2.2pt);
\fill[navy] (2,6) circle (2.2pt);
\fill[navy] (4,14) circle (2.2pt);
\fill[navy] (5,14) circle (2.2pt);
\fill[navy] (7,0) circle (2.2pt);
\end{tikzpicture}\end{center}\begin{tasks}(1)
  \task A quina distància de casa estava a les 3 hores? I a les 6 hores?
  \task Quantes vegades va parar i quant de temps cada vegada?
  \task Quina és la distància màxima a què va arribar? A quina hora?
  \task Quant va tardar a tornar? Quina distància va recórrer en total?
\end{tasks}\recorda{Gràfic de distància i temps: si puja, s'allunya; si és horitzontal, està parat; si baixa, torna. Com més inclinat, més de pressa.}

\ex{1} El gràfic dona la temperatura al llarg d'un dia.\par\smallskip\begin{center}\begin{tikzpicture}[x=0.26cm,y=0.2cm,baseline=(current bounding box.center),line cap=round,line join=round]
\draw[muted!30,line width=0.3pt] (0,0)--(0,20);
\draw[muted!30,line width=0.3pt] (2,0)--(2,20);
\draw[muted!30,line width=0.3pt] (4,0)--(4,20);
\draw[muted!30,line width=0.3pt] (6,0)--(6,20);
\draw[muted!30,line width=0.3pt] (8,0)--(8,20);
\draw[muted!30,line width=0.3pt] (10,0)--(10,20);
\draw[muted!30,line width=0.3pt] (12,0)--(12,20);
\draw[muted!30,line width=0.3pt] (14,0)--(14,20);
\draw[muted!30,line width=0.3pt] (16,0)--(16,20);
\draw[muted!30,line width=0.3pt] (18,0)--(18,20);
\draw[muted!30,line width=0.3pt] (20,0)--(20,20);
\draw[muted!30,line width=0.3pt] (22,0)--(22,20);
\draw[muted!30,line width=0.3pt] (24,0)--(24,20);
\draw[muted!30,line width=0.3pt] (0,0)--(25,0);
\draw[muted!30,line width=0.3pt] (0,2)--(25,2);
\draw[muted!30,line width=0.3pt] (0,4)--(25,4);
\draw[muted!30,line width=0.3pt] (0,6)--(25,6);
\draw[muted!30,line width=0.3pt] (0,8)--(25,8);
\draw[muted!30,line width=0.3pt] (0,10)--(25,10);
\draw[muted!30,line width=0.3pt] (0,12)--(25,12);
\draw[muted!30,line width=0.3pt] (0,14)--(25,14);
\draw[muted!30,line width=0.3pt] (0,16)--(25,16);
\draw[muted!30,line width=0.3pt] (0,18)--(25,18);
\draw[muted!30,line width=0.3pt] (0,20)--(25,20);
\draw[-latex,muted,line width=0.6pt] (0,0)--(25.6,0) node[right,font=\scriptsize,text=navy]{h};
\draw[-latex,muted,line width=0.6pt] (0,0)--(0,20.6) node[above,font=\scriptsize,text=navy]{$^\circ$C};
\node[font=\tiny,text=muted,below,inner sep=1.5pt] at (2,0) {$2$};
\node[font=\tiny,text=muted,below,inner sep=1.5pt] at (4,0) {$4$};
\node[font=\tiny,text=muted,below,inner sep=1.5pt] at (6,0) {$6$};
\node[font=\tiny,text=muted,below,inner sep=1.5pt] at (8,0) {$8$};
\node[font=\tiny,text=muted,below,inner sep=1.5pt] at (10,0) {$10$};
\node[font=\tiny,text=muted,below,inner sep=1.5pt] at (12,0) {$12$};
\node[font=\tiny,text=muted,below,inner sep=1.5pt] at (14,0) {$14$};
\node[font=\tiny,text=muted,below,inner sep=1.5pt] at (16,0) {$16$};
\node[font=\tiny,text=muted,below,inner sep=1.5pt] at (18,0) {$18$};
\node[font=\tiny,text=muted,below,inner sep=1.5pt] at (20,0) {$20$};
\node[font=\tiny,text=muted,below,inner sep=1.5pt] at (22,0) {$22$};
\node[font=\tiny,text=muted,below,inner sep=1.5pt] at (24,0) {$24$};
\node[font=\tiny,text=muted,left,inner sep=1.5pt] at (0,2) {$2$};
\node[font=\tiny,text=muted,left,inner sep=1.5pt] at (0,4) {$4$};
\node[font=\tiny,text=muted,left,inner sep=1.5pt] at (0,6) {$6$};
\node[font=\tiny,text=muted,left,inner sep=1.5pt] at (0,8) {$8$};
\node[font=\tiny,text=muted,left,inner sep=1.5pt] at (0,10) {$10$};
\node[font=\tiny,text=muted,left,inner sep=1.5pt] at (0,12) {$12$};
\node[font=\tiny,text=muted,left,inner sep=1.5pt] at (0,14) {$14$};
\node[font=\tiny,text=muted,left,inner sep=1.5pt] at (0,16) {$16$};
\node[font=\tiny,text=muted,left,inner sep=1.5pt] at (0,18) {$18$};
\node[font=\tiny,text=muted,below left,inner sep=1.5pt] at (0,0) {$0$};
\draw[accent,line width=1.1pt] (0,6)--(2,4)--(6,4)--(10,12)--(14,18)--(18,14)--(22,8)--(24,7);
\end{tikzpicture}\end{center}\begin{tasks}(1)
  \task Quina va ser la temperatura màxima? I la mínima?
  \task En quins intervals de temps la temperatura puja? En quins baixa? Quan es manté constant?
\end{tasks}

\ex{1} El gràfic mostra l'altura d'una planta durant 6 setmanes.\par\smallskip\begin{center}\begin{tikzpicture}[x=0.6cm,y=0.22cm,baseline=(current bounding box.center),line cap=round,line join=round]
\draw[muted!30,line width=0.3pt] (0,0)--(0,16.5);
\draw[muted!30,line width=0.3pt] (1,0)--(1,16.5);
\draw[muted!30,line width=0.3pt] (2,0)--(2,16.5);
\draw[muted!30,line width=0.3pt] (3,0)--(3,16.5);
\draw[muted!30,line width=0.3pt] (4,0)--(4,16.5);
\draw[muted!30,line width=0.3pt] (5,0)--(5,16.5);
\draw[muted!30,line width=0.3pt] (6,0)--(6,16.5);
\draw[muted!30,line width=0.3pt] (0,0)--(6.6,0);
\draw[muted!30,line width=0.3pt] (0,2)--(6.6,2);
\draw[muted!30,line width=0.3pt] (0,4)--(6.6,4);
\draw[muted!30,line width=0.3pt] (0,6)--(6.6,6);
\draw[muted!30,line width=0.3pt] (0,8)--(6.6,8);
\draw[muted!30,line width=0.3pt] (0,10)--(6.6,10);
\draw[muted!30,line width=0.3pt] (0,12)--(6.6,12);
\draw[muted!30,line width=0.3pt] (0,14)--(6.6,14);
\draw[muted!30,line width=0.3pt] (0,16)--(6.6,16);
\draw[-latex,muted,line width=0.6pt] (0,0)--(6.9,0) node[right,font=\scriptsize,text=navy]{setmanes};
\draw[-latex,muted,line width=0.6pt] (0,0)--(0,17.1) node[above,font=\scriptsize,text=navy]{cm};
\node[font=\tiny,text=muted,below,inner sep=1.5pt] at (1,0) {$1$};
\node[font=\tiny,text=muted,below,inner sep=1.5pt] at (2,0) {$2$};
\node[font=\tiny,text=muted,below,inner sep=1.5pt] at (3,0) {$3$};
\node[font=\tiny,text=muted,below,inner sep=1.5pt] at (4,0) {$4$};
\node[font=\tiny,text=muted,below,inner sep=1.5pt] at (5,0) {$5$};
\node[font=\tiny,text=muted,below,inner sep=1.5pt] at (6,0) {$6$};
\node[font=\tiny,text=muted,left,inner sep=1.5pt] at (0,2) {$2$};
\node[font=\tiny,text=muted,left,inner sep=1.5pt] at (0,4) {$4$};
\node[font=\tiny,text=muted,left,inner sep=1.5pt] at (0,6) {$6$};
\node[font=\tiny,text=muted,left,inner sep=1.5pt] at (0,8) {$8$};
\node[font=\tiny,text=muted,left,inner sep=1.5pt] at (0,10) {$10$};
\node[font=\tiny,text=muted,left,inner sep=1.5pt] at (0,12) {$12$};
\node[font=\tiny,text=muted,left,inner sep=1.5pt] at (0,14) {$14$};
\node[font=\tiny,text=muted,left,inner sep=1.5pt] at (0,16) {$16$};
\node[font=\tiny,text=muted,below left,inner sep=1.5pt] at (0,0) {$0$};
\draw[teal,line width=1.1pt] (0,2)--(1,5)--(2,9)--(3,12)--(4,14)--(5,15)--(6,15);
\fill[navy] (0,2) circle (2.2pt);
\fill[navy] (1,5) circle (2.2pt);
\fill[navy] (2,9) circle (2.2pt);
\fill[navy] (3,12) circle (2.2pt);
\fill[navy] (4,14) circle (2.2pt);
\fill[navy] (5,15) circle (2.2pt);
\fill[navy] (6,15) circle (2.2pt);
\end{tikzpicture}\end{center}\begin{tasks}(1)
  \task Quant feia en començar? I a la 4a setmana?
  \task En quina setmana va créixer més?
  \task Què passa entre la 5a i la 6a setmana?
\end{tasks}

\ex{2} El gràfic mostra l'altura de l'aigua d'una banyera. Explica què passa en cada tram.\par\smallskip\begin{center}\begin{tikzpicture}[x=0.33cm,y=0.14cm,baseline=(current bounding box.center),line cap=round,line join=round]
\draw[muted!30,line width=0.3pt] (0,0)--(0,30);
\draw[muted!30,line width=0.3pt] (2,0)--(2,30);
\draw[muted!30,line width=0.3pt] (4,0)--(4,30);
\draw[muted!30,line width=0.3pt] (6,0)--(6,30);
\draw[muted!30,line width=0.3pt] (8,0)--(8,30);
\draw[muted!30,line width=0.3pt] (10,0)--(10,30);
\draw[muted!30,line width=0.3pt] (12,0)--(12,30);
\draw[muted!30,line width=0.3pt] (14,0)--(14,30);
\draw[muted!30,line width=0.3pt] (16,0)--(16,30);
\draw[muted!30,line width=0.3pt] (18,0)--(18,30);
\draw[muted!30,line width=0.3pt] (20,0)--(20,30);
\draw[muted!30,line width=0.3pt] (22,0)--(22,30);
\draw[muted!30,line width=0.3pt] (0,0)--(23,0);
\draw[muted!30,line width=0.3pt] (0,5)--(23,5);
\draw[muted!30,line width=0.3pt] (0,10)--(23,10);
\draw[muted!30,line width=0.3pt] (0,15)--(23,15);
\draw[muted!30,line width=0.3pt] (0,20)--(23,20);
\draw[muted!30,line width=0.3pt] (0,25)--(23,25);
\draw[muted!30,line width=0.3pt] (0,30)--(23,30);
\draw[-latex,muted,line width=0.6pt] (0,0)--(23.6,0) node[right,font=\scriptsize,text=navy]{min};
\draw[-latex,muted,line width=0.6pt] (0,0)--(0,31.5) node[above,font=\scriptsize,text=navy]{cm};
\node[font=\tiny,text=muted,below,inner sep=1.5pt] at (2,0) {$2$};
\node[font=\tiny,text=muted,below,inner sep=1.5pt] at (4,0) {$4$};
\node[font=\tiny,text=muted,below,inner sep=1.5pt] at (6,0) {$6$};
\node[font=\tiny,text=muted,below,inner sep=1.5pt] at (8,0) {$8$};
\node[font=\tiny,text=muted,below,inner sep=1.5pt] at (10,0) {$10$};
\node[font=\tiny,text=muted,below,inner sep=1.5pt] at (12,0) {$12$};
\node[font=\tiny,text=muted,below,inner sep=1.5pt] at (14,0) {$14$};
\node[font=\tiny,text=muted,below,inner sep=1.5pt] at (16,0) {$16$};
\node[font=\tiny,text=muted,below,inner sep=1.5pt] at (18,0) {$18$};
\node[font=\tiny,text=muted,below,inner sep=1.5pt] at (20,0) {$20$};
\node[font=\tiny,text=muted,below,inner sep=1.5pt] at (22,0) {$22$};
\node[font=\tiny,text=muted,left,inner sep=1.5pt] at (0,5) {$5$};
\node[font=\tiny,text=muted,left,inner sep=1.5pt] at (0,10) {$10$};
\node[font=\tiny,text=muted,left,inner sep=1.5pt] at (0,15) {$15$};
\node[font=\tiny,text=muted,left,inner sep=1.5pt] at (0,20) {$20$};
\node[font=\tiny,text=muted,left,inner sep=1.5pt] at (0,25) {$25$};
\node[font=\tiny,text=muted,below left,inner sep=1.5pt] at (0,0) {$0$};
\draw[accent,line width=1.1pt] (0,0)--(5,20)--(6,26)--(16,26)--(17,20)--(22,0);
\end{tikzpicture}\end{center}

\ex{2} Quin gràfic (temps, distància a casa) correspon a cada història? N'hi ha un que sobra.\par\smallskip\begin{center}\begin{tabular}{c}\begin{tikzpicture}[x=0.3cm,y=0.22cm,baseline=(current bounding box.center),line cap=round,line join=round]
\draw[muted!30,line width=0.3pt] (0,0)--(0,9);
\draw[muted!30,line width=0.3pt] (2,0)--(2,9);
\draw[muted!30,line width=0.3pt] (4,0)--(4,9);
\draw[muted!30,line width=0.3pt] (6,0)--(6,9);
\draw[muted!30,line width=0.3pt] (0,0)--(6.5,0);
\draw[muted!30,line width=0.3pt] (0,2)--(6.5,2);
\draw[muted!30,line width=0.3pt] (0,4)--(6.5,4);
\draw[muted!30,line width=0.3pt] (0,6)--(6.5,6);
\draw[muted!30,line width=0.3pt] (0,8)--(6.5,8);
\draw[-latex,muted,line width=0.6pt] (0,0)--(7.1,0) node[right,font=\scriptsize,text=navy]{$t$};
\draw[-latex,muted,line width=0.6pt] (0,0)--(0,9.6) node[above,font=\scriptsize,text=navy]{$d$};
\node[font=\tiny,text=muted,below,inner sep=1.5pt] at (2,0) {$2$};
\node[font=\tiny,text=muted,below,inner sep=1.5pt] at (4,0) {$4$};
\node[font=\tiny,text=muted,below,inner sep=1.5pt] at (6,0) {$6$};
\node[font=\tiny,text=muted,left,inner sep=1.5pt] at (0,2) {$2$};
\node[font=\tiny,text=muted,left,inner sep=1.5pt] at (0,4) {$4$};
\node[font=\tiny,text=muted,left,inner sep=1.5pt] at (0,6) {$6$};
\node[font=\tiny,text=muted,left,inner sep=1.5pt] at (0,8) {$8$};
\node[font=\tiny,text=muted,below left,inner sep=1.5pt] at (0,0) {$0$};
\draw[accent,line width=1.1pt] (0,0)--(2,4)--(4,4)--(6,0);
\end{tikzpicture}\\ \textbf{A}\end{tabular}\hspace{0.3cm}\begin{tabular}{c}\begin{tikzpicture}[x=0.3cm,y=0.22cm,baseline=(current bounding box.center),line cap=round,line join=round]
\draw[muted!30,line width=0.3pt] (0,0)--(0,9);
\draw[muted!30,line width=0.3pt] (2,0)--(2,9);
\draw[muted!30,line width=0.3pt] (4,0)--(4,9);
\draw[muted!30,line width=0.3pt] (6,0)--(6,9);
\draw[muted!30,line width=0.3pt] (0,0)--(6.5,0);
\draw[muted!30,line width=0.3pt] (0,2)--(6.5,2);
\draw[muted!30,line width=0.3pt] (0,4)--(6.5,4);
\draw[muted!30,line width=0.3pt] (0,6)--(6.5,6);
\draw[muted!30,line width=0.3pt] (0,8)--(6.5,8);
\draw[-latex,muted,line width=0.6pt] (0,0)--(7.1,0) node[right,font=\scriptsize,text=navy]{$t$};
\draw[-latex,muted,line width=0.6pt] (0,0)--(0,9.6) node[above,font=\scriptsize,text=navy]{$d$};
\node[font=\tiny,text=muted,below,inner sep=1.5pt] at (2,0) {$2$};
\node[font=\tiny,text=muted,below,inner sep=1.5pt] at (4,0) {$4$};
\node[font=\tiny,text=muted,below,inner sep=1.5pt] at (6,0) {$6$};
\node[font=\tiny,text=muted,left,inner sep=1.5pt] at (0,2) {$2$};
\node[font=\tiny,text=muted,left,inner sep=1.5pt] at (0,4) {$4$};
\node[font=\tiny,text=muted,left,inner sep=1.5pt] at (0,6) {$6$};
\node[font=\tiny,text=muted,left,inner sep=1.5pt] at (0,8) {$8$};
\node[font=\tiny,text=muted,below left,inner sep=1.5pt] at (0,0) {$0$};
\draw[accent,line width=1.1pt] (0,0)--(6,8);
\end{tikzpicture}\\ \textbf{B}\end{tabular}\hspace{0.3cm}\begin{tabular}{c}\begin{tikzpicture}[x=0.3cm,y=0.22cm,baseline=(current bounding box.center),line cap=round,line join=round]
\draw[muted!30,line width=0.3pt] (0,0)--(0,9);
\draw[muted!30,line width=0.3pt] (2,0)--(2,9);
\draw[muted!30,line width=0.3pt] (4,0)--(4,9);
\draw[muted!30,line width=0.3pt] (6,0)--(6,9);
\draw[muted!30,line width=0.3pt] (0,0)--(6.5,0);
\draw[muted!30,line width=0.3pt] (0,2)--(6.5,2);
\draw[muted!30,line width=0.3pt] (0,4)--(6.5,4);
\draw[muted!30,line width=0.3pt] (0,6)--(6.5,6);
\draw[muted!30,line width=0.3pt] (0,8)--(6.5,8);
\draw[-latex,muted,line width=0.6pt] (0,0)--(7.1,0) node[right,font=\scriptsize,text=navy]{$t$};
\draw[-latex,muted,line width=0.6pt] (0,0)--(0,9.6) node[above,font=\scriptsize,text=navy]{$d$};
\node[font=\tiny,text=muted,below,inner sep=1.5pt] at (2,0) {$2$};
\node[font=\tiny,text=muted,below,inner sep=1.5pt] at (4,0) {$4$};
\node[font=\tiny,text=muted,below,inner sep=1.5pt] at (6,0) {$6$};
\node[font=\tiny,text=muted,left,inner sep=1.5pt] at (0,2) {$2$};
\node[font=\tiny,text=muted,left,inner sep=1.5pt] at (0,4) {$4$};
\node[font=\tiny,text=muted,left,inner sep=1.5pt] at (0,6) {$6$};
\node[font=\tiny,text=muted,left,inner sep=1.5pt] at (0,8) {$8$};
\node[font=\tiny,text=muted,below left,inner sep=1.5pt] at (0,0) {$0$};
\draw[accent,line width=1.1pt] (0,0)--(2,4)--(4,4)--(6,8);
\end{tikzpicture}\\ \textbf{C}\end{tabular}\hspace{0.3cm}\begin{tabular}{c}\begin{tikzpicture}[x=0.3cm,y=0.22cm,baseline=(current bounding box.center),line cap=round,line join=round]
\draw[muted!30,line width=0.3pt] (0,0)--(0,9);
\draw[muted!30,line width=0.3pt] (2,0)--(2,9);
\draw[muted!30,line width=0.3pt] (4,0)--(4,9);
\draw[muted!30,line width=0.3pt] (6,0)--(6,9);
\draw[muted!30,line width=0.3pt] (0,0)--(6.5,0);
\draw[muted!30,line width=0.3pt] (0,2)--(6.5,2);
\draw[muted!30,line width=0.3pt] (0,4)--(6.5,4);
\draw[muted!30,line width=0.3pt] (0,6)--(6.5,6);
\draw[muted!30,line width=0.3pt] (0,8)--(6.5,8);
\draw[-latex,muted,line width=0.6pt] (0,0)--(7.1,0) node[right,font=\scriptsize,text=navy]{$t$};
\draw[-latex,muted,line width=0.6pt] (0,0)--(0,9.6) node[above,font=\scriptsize,text=navy]{$d$};
\node[font=\tiny,text=muted,below,inner sep=1.5pt] at (2,0) {$2$};
\node[font=\tiny,text=muted,below,inner sep=1.5pt] at (4,0) {$4$};
\node[font=\tiny,text=muted,below,inner sep=1.5pt] at (6,0) {$6$};
\node[font=\tiny,text=muted,left,inner sep=1.5pt] at (0,2) {$2$};
\node[font=\tiny,text=muted,left,inner sep=1.5pt] at (0,4) {$4$};
\node[font=\tiny,text=muted,left,inner sep=1.5pt] at (0,6) {$6$};
\node[font=\tiny,text=muted,left,inner sep=1.5pt] at (0,8) {$8$};
\node[font=\tiny,text=muted,below left,inner sep=1.5pt] at (0,0) {$0$};
\draw[accent,line width=1.1pt] (0,8)--(3,2)--(6,2);
\end{tikzpicture}\\ \textbf{D}\end{tabular}\end{center}\begin{tasks}(1)
  \task La Sara va a casa d'una amiga sense parar.
  \task L'Òscar ix de casa, para a comprar el pa i continua fins a l'escola.
  \task La Lucía va al quiosc, s'hi queda una estona i torna a casa.
\end{tasks}

\ex{2} El gràfic mostra el viatge d'un tren.\par\smallskip\begin{center}\begin{tikzpicture}[x=0.12cm,y=0.05cm,baseline=(current bounding box.center),line cap=round,line join=round]
\draw[muted!30,line width=0.3pt] (0,0)--(0,100);
\draw[muted!30,line width=0.3pt] (5,0)--(5,100);
\draw[muted!30,line width=0.3pt] (10,0)--(10,100);
\draw[muted!30,line width=0.3pt] (15,0)--(15,100);
\draw[muted!30,line width=0.3pt] (20,0)--(20,100);
\draw[muted!30,line width=0.3pt] (25,0)--(25,100);
\draw[muted!30,line width=0.3pt] (30,0)--(30,100);
\draw[muted!30,line width=0.3pt] (35,0)--(35,100);
\draw[muted!30,line width=0.3pt] (40,0)--(40,100);
\draw[muted!30,line width=0.3pt] (45,0)--(45,100);
\draw[muted!30,line width=0.3pt] (50,0)--(50,100);
\draw[muted!30,line width=0.3pt] (55,0)--(55,100);
\draw[muted!30,line width=0.3pt] (60,0)--(60,100);
\draw[muted!30,line width=0.3pt] (65,0)--(65,100);
\draw[muted!30,line width=0.3pt] (70,0)--(70,100);
\draw[muted!30,line width=0.3pt] (75,0)--(75,100);
\draw[muted!30,line width=0.3pt] (0,0)--(75,0);
\draw[muted!30,line width=0.3pt] (0,10)--(75,10);
\draw[muted!30,line width=0.3pt] (0,20)--(75,20);
\draw[muted!30,line width=0.3pt] (0,30)--(75,30);
\draw[muted!30,line width=0.3pt] (0,40)--(75,40);
\draw[muted!30,line width=0.3pt] (0,50)--(75,50);
\draw[muted!30,line width=0.3pt] (0,60)--(75,60);
\draw[muted!30,line width=0.3pt] (0,70)--(75,70);
\draw[muted!30,line width=0.3pt] (0,80)--(75,80);
\draw[muted!30,line width=0.3pt] (0,90)--(75,90);
\draw[muted!30,line width=0.3pt] (0,100)--(75,100);
\draw[-latex,muted,line width=0.6pt] (0,0)--(76.5,0) node[right,font=\scriptsize,text=navy]{min};
\draw[-latex,muted,line width=0.6pt] (0,0)--(0,103) node[above,font=\scriptsize,text=navy]{km};
\node[font=\tiny,text=muted,below,inner sep=1.5pt] at (10,0) {$10$};
\node[font=\tiny,text=muted,below,inner sep=1.5pt] at (20,0) {$20$};
\node[font=\tiny,text=muted,below,inner sep=1.5pt] at (30,0) {$30$};
\node[font=\tiny,text=muted,below,inner sep=1.5pt] at (40,0) {$40$};
\node[font=\tiny,text=muted,below,inner sep=1.5pt] at (50,0) {$50$};
\node[font=\tiny,text=muted,below,inner sep=1.5pt] at (60,0) {$60$};
\node[font=\tiny,text=muted,below,inner sep=1.5pt] at (70,0) {$70$};
\node[font=\tiny,text=muted,left,inner sep=1.5pt] at (0,20) {$20$};
\node[font=\tiny,text=muted,left,inner sep=1.5pt] at (0,40) {$40$};
\node[font=\tiny,text=muted,left,inner sep=1.5pt] at (0,60) {$60$};
\node[font=\tiny,text=muted,left,inner sep=1.5pt] at (0,80) {$80$};
\node[font=\tiny,text=muted,below left,inner sep=1.5pt] at (0,0) {$0$};
\draw[accent,line width=1.1pt] (0,0)--(20,30)--(25,30)--(45,60)--(50,60)--(70,90);
\end{tikzpicture}\end{center}\begin{tasks}(1)
  \task Quantes parades fa i quant duren?
  \task Quants km fa entre parada i parada? Quant tarda?
  \task Quina distància total recorre i quant de temps tarda?
\end{tasks}

\ex{3} L'Anna i el Pau fan la mateixa ruta en bici. Observa el gràfic i respon.\par\smallskip\begin{center}\begin{tikzpicture}[x=0.9cm,y=0.04cm,baseline=(current bounding box.center),line cap=round,line join=round]
\draw[muted!30,line width=0.3pt] (0,0)--(0,105);
\draw[muted!30,line width=0.3pt] (1,0)--(1,105);
\draw[muted!30,line width=0.3pt] (2,0)--(2,105);
\draw[muted!30,line width=0.3pt] (3,0)--(3,105);
\draw[muted!30,line width=0.3pt] (4,0)--(4,105);
\draw[muted!30,line width=0.3pt] (5,0)--(5,105);
\draw[muted!30,line width=0.3pt] (0,0)--(5.5,0);
\draw[muted!30,line width=0.3pt] (0,10)--(5.5,10);
\draw[muted!30,line width=0.3pt] (0,20)--(5.5,20);
\draw[muted!30,line width=0.3pt] (0,30)--(5.5,30);
\draw[muted!30,line width=0.3pt] (0,40)--(5.5,40);
\draw[muted!30,line width=0.3pt] (0,50)--(5.5,50);
\draw[muted!30,line width=0.3pt] (0,60)--(5.5,60);
\draw[muted!30,line width=0.3pt] (0,70)--(5.5,70);
\draw[muted!30,line width=0.3pt] (0,80)--(5.5,80);
\draw[muted!30,line width=0.3pt] (0,90)--(5.5,90);
\draw[muted!30,line width=0.3pt] (0,100)--(5.5,100);
\draw[-latex,muted,line width=0.6pt] (0,0)--(5.8,0) node[right,font=\scriptsize,text=navy]{h};
\draw[-latex,muted,line width=0.6pt] (0,0)--(0,108) node[above,font=\scriptsize,text=navy]{km};
\node[font=\tiny,text=muted,below,inner sep=1.5pt] at (1,0) {$1$};
\node[font=\tiny,text=muted,below,inner sep=1.5pt] at (2,0) {$2$};
\node[font=\tiny,text=muted,below,inner sep=1.5pt] at (3,0) {$3$};
\node[font=\tiny,text=muted,below,inner sep=1.5pt] at (4,0) {$4$};
\node[font=\tiny,text=muted,below,inner sep=1.5pt] at (5,0) {$5$};
\node[font=\tiny,text=muted,left,inner sep=1.5pt] at (0,20) {$20$};
\node[font=\tiny,text=muted,left,inner sep=1.5pt] at (0,40) {$40$};
\node[font=\tiny,text=muted,left,inner sep=1.5pt] at (0,60) {$60$};
\node[font=\tiny,text=muted,left,inner sep=1.5pt] at (0,80) {$80$};
\node[font=\tiny,text=muted,left,inner sep=1.5pt] at (0,100) {$100$};
\node[font=\tiny,text=muted,below left,inner sep=1.5pt] at (0,0) {$0$};
\draw[accent,line width=1.1pt] (0,0)--(5,75);
\draw[teal,line width=1.1pt] (1,0)--(5,100);
\fill[navy] (2.5,37.5) circle (2.2pt);
\node[font=\scriptsize,right,text=accent] at (5,75) {Anna};
\node[font=\scriptsize,left,text=teal] at (4.6,100) {Pau};
\end{tikzpicture}\end{center}\begin{tasks}(1)
  \task Qui ix primer? Quanta estona abans?
  \task Quants km fa cadascun en una hora? Qui va més de pressa?
  \task On i quan atrapa el Pau l'Anna?
\end{tasks}

\ex{3} En Marc ix de casa, camina 1\,km en 20 minuts fins a la biblioteca, hi està mitja hora i torna a casa en 10 minuts. Dibuixa el gràfic (temps en minuts, distància a casa en km). En quin tram va més de pressa?


% =====================================================
\seccio{5}{Funció de proporcionalitat directa: $y=m\cdot x$}

\ex{1} Quines d'aquestes funcions són de proporcionalitat directa? Escriu-ne el pendent $m$.\begin{tasks}(6)
  \task $y=4x$
  \task $y=x+4$
  \task $y=-\tfrac{x}{2}$
  \task $y=4$
  \task $y=0{,}6x$
  \task $y=-3x$
\end{tasks}\recorda{Una funció de proporcionalitat directa és de la forma $y=m\cdot x$. El seu gràfic és una recta que passa per l'origen. El pendent $m$ és el que augmenta $y$ quan $x$ augmenta 1.}

\ex{1} Completa la taula de $y=3x$ i dibuixa'n el gràfic.\par\smallskip\begin{center}$\begin{array}{|c|C{0.75cm}C{0.75cm}C{0.75cm}C{0.75cm}C{0.75cm}|}\hline x&-2&-1&0&1&2\\\hline y&\phantom{00}&\phantom{00}&\phantom{00}&\phantom{00}&\phantom{00}\\\hline\end{array}$\end{center}

\ex{1} Dibuixa en els mateixos eixos $y=-2x$ i $y=\dfrac{x}{2}$. Quina és creixent? Quina és més inclinada?

\ex{2} Una recta passa per l'origen i pel punt $(3,\ -6)$. Escriu-ne la fórmula i comprova si els punts $(-4,\ 8)$ i $(5,\ -12)$ són de la recta.

\ex{2} Comprova que aquesta taula és de proporcionalitat directa, escriu-ne la fórmula i dibuixa el gràfic. Quant val $y$ si $x=10$?\par\smallskip\begin{center}$\begin{array}{|c|cccc|}\hline x&2&4&6&8\\\hline y&3&6&9&12\\\hline\end{array}$\end{center}

\ex{2} Ordena aquestes rectes de més inclinada a menys i digues quines són decreixents: \ $a$: $y=5x$, \ $b$: $y=-\frac12x$, \ $c$: $y=\frac13x$, \ $d$: $y=-2x$.

\ex{3} Un dia, 1\,€ val 1,10 dòlars. Escriu la fórmula que passa d'euros ($x$) a dòlars ($y$) i dibuixa'n el gràfic. Quants dòlars són 20\,€? I 35\,€? Quants euros són 55\,\$?

\ex{3} Un cotxe va a 90\,km/h sempre a la mateixa velocitat. Escriu la fórmula de la distància $d$ que recorre en $t$ hores, dibuixa'n el gràfic i calcula quina distància fa en 3 hores i quant tarda a fer 225\,km.


% =====================================================
\seccio{6}{Funció lineal: $y=m\cdot x+n$}

\ex{1} Escriu el pendent $m$ i l'ordenada en l'origen $n$ de cada funció:\begin{tasks}(6)
  \task $y=3x-4$
  \task $y=-x+5$
  \task $y=\tfrac{x}{2}+1$
  \task $y=-2x$
  \task $y=7$
  \task $y=6-4x$
\end{tasks}\recorda{En $y=mx+n$, $m$ és el \textbf{pendent} (el que canvia $y$ quan $x$ augmenta 1) i $n$ és l'\textbf{ordenada en l'origen}: la recta talla l'eix $Y$ en $(0,\ n)$.}

\ex{1} Completa la taula de $y=2x-1$ i dibuixa'n el gràfic. On talla l'eix $Y$?\par\smallskip\begin{center}$\begin{array}{|c|C{0.75cm}C{0.75cm}C{0.75cm}C{0.75cm}C{0.75cm}|}\hline x&-1&0&1&2&3\\\hline y&\phantom{00}&\phantom{00}&\phantom{00}&\phantom{00}&\phantom{00}\\\hline\end{array}$\end{center}

\ex{1} Quins d'aquests punts són de la recta $y=-2x+3$? \ $(2,\ 1)$, \ $(-1,\ 7)$, \ $(3,\ -2)$, \ $(0,\ 3)$.

\ex{2} Dibuixa en els mateixos eixos $y=x+2$ i $y=-x+4$. En quin punt es tallen? Comprova-ho amb les fórmules.

\ex{2} Escriu la fórmula de cadascuna de les rectes del dibuix.\par\smallskip\begin{center}\begin{tikzpicture}[x=0.4cm,y=0.4cm,baseline=(current bounding box.center),line cap=round,line join=round]
\draw[muted!30,line width=0.3pt] (-5,-5)--(-5,5);
\draw[muted!30,line width=0.3pt] (-4,-5)--(-4,5);
\draw[muted!30,line width=0.3pt] (-3,-5)--(-3,5);
\draw[muted!30,line width=0.3pt] (-2,-5)--(-2,5);
\draw[muted!30,line width=0.3pt] (-1,-5)--(-1,5);
\draw[muted!30,line width=0.3pt] (0,-5)--(0,5);
\draw[muted!30,line width=0.3pt] (1,-5)--(1,5);
\draw[muted!30,line width=0.3pt] (2,-5)--(2,5);
\draw[muted!30,line width=0.3pt] (3,-5)--(3,5);
\draw[muted!30,line width=0.3pt] (4,-5)--(4,5);
\draw[muted!30,line width=0.3pt] (5,-5)--(5,5);
\draw[muted!30,line width=0.3pt] (-5,-5)--(5,-5);
\draw[muted!30,line width=0.3pt] (-5,-4)--(5,-4);
\draw[muted!30,line width=0.3pt] (-5,-3)--(5,-3);
\draw[muted!30,line width=0.3pt] (-5,-2)--(5,-2);
\draw[muted!30,line width=0.3pt] (-5,-1)--(5,-1);
\draw[muted!30,line width=0.3pt] (-5,0)--(5,0);
\draw[muted!30,line width=0.3pt] (-5,1)--(5,1);
\draw[muted!30,line width=0.3pt] (-5,2)--(5,2);
\draw[muted!30,line width=0.3pt] (-5,3)--(5,3);
\draw[muted!30,line width=0.3pt] (-5,4)--(5,4);
\draw[muted!30,line width=0.3pt] (-5,5)--(5,5);
\draw[-latex,muted,line width=0.6pt] (-5,0)--(5.3,0) node[right,font=\scriptsize,text=navy]{x};
\draw[-latex,muted,line width=0.6pt] (0,-5)--(0,5.3) node[above,font=\scriptsize,text=navy]{y};
\node[font=\tiny,text=muted,below,inner sep=1.5pt] at (-4,0) {$-4$};
\node[font=\tiny,text=muted,below,inner sep=1.5pt] at (-3,0) {$-3$};
\node[font=\tiny,text=muted,below,inner sep=1.5pt] at (-2,0) {$-2$};
\node[font=\tiny,text=muted,below,inner sep=1.5pt] at (-1,0) {$-1$};
\node[font=\tiny,text=muted,below,inner sep=1.5pt] at (1,0) {$1$};
\node[font=\tiny,text=muted,below,inner sep=1.5pt] at (2,0) {$2$};
\node[font=\tiny,text=muted,below,inner sep=1.5pt] at (3,0) {$3$};
\node[font=\tiny,text=muted,below,inner sep=1.5pt] at (4,0) {$4$};
\node[font=\tiny,text=muted,left,inner sep=1.5pt] at (0,-4) {$-4$};
\node[font=\tiny,text=muted,left,inner sep=1.5pt] at (0,-3) {$-3$};
\node[font=\tiny,text=muted,left,inner sep=1.5pt] at (0,-2) {$-2$};
\node[font=\tiny,text=muted,left,inner sep=1.5pt] at (0,-1) {$-1$};
\node[font=\tiny,text=muted,left,inner sep=1.5pt] at (0,1) {$1$};
\node[font=\tiny,text=muted,left,inner sep=1.5pt] at (0,2) {$2$};
\node[font=\tiny,text=muted,left,inner sep=1.5pt] at (0,3) {$3$};
\node[font=\tiny,text=muted,left,inner sep=1.5pt] at (0,4) {$4$};
\node[font=\tiny,text=muted,below left,inner sep=1.5pt] at (0,0) {$0$};
\draw[accent,line width=1.1pt] (-2.000,-5.000)--(3.000,5.000);
\draw[teal,line width=1.1pt] (-2.000,5.000)--(5.000,-2.000);
\draw[navy,line width=1.1pt] (-5.000,-1.500)--(5.000,3.500);
\draw[violet,line width=1.1pt] (-5.000,-3.000)--(5.000,-3.000);
\node[font=\scriptsize,right,text=accent] at (2.8,4.6) {$a$};
\node[font=\scriptsize,left,text=teal] at (-1.8,4.6) {$b$};
\node[font=\scriptsize,above,text=navy] at (4.4,3.2) {$c$};
\node[font=\scriptsize,above,text=violet] at (-4.5,-3) {$d$};
\end{tikzpicture}\end{center}

\ex{2} Observa aquestes funcions: \ $y=3x+1$, \ $y=-x$, \ $y=3x-4$, \ $y=5$, \ $y=-x+2$, \ $y=2$. Quines són paral·leles? Quines passen per l'origen? Quines són constants?

\ex{2} Troba la fórmula de la funció lineal d'aquesta taula:\par\smallskip\begin{center}$\begin{array}{|c|cccc|}\hline x&0&1&2&3\\\hline y&4&1&-2&-5\\\hline\end{array}$\end{center}

\ex{3} Troba la fórmula de la recta que passa per:\begin{tasks}(2)
  \task $(0,\ 3)$ i $(2,\ 7)$
  \task $(1,\ -1)$ i $(3,\ 3)$
\end{tasks}

\ex{3} Respon:\begin{tasks}(1)
  \task Quin valor ha de tindre $k$ perquè el punt $(2,\ k)$ siga de la recta $y=-3x+5$?
  \task Quin valor ha de tindre $a$ perquè el punt $(a,\ 11)$ siga de la recta $y=2x+5$?
\end{tasks}


% =====================================================
\seccio{7}{Problemes amb funcions lineals}

\ex{1} Un taxi cobra 3,50\,€ de baixada de bandera i 1,50\,€ per km. Escriu la fórmula del preu, dibuixa'n el gràfic i calcula quant costa un viatge de 8\,km. Quants km hem fet si paguem 26\,€?\recorda{En molts problemes hi ha una \textbf{part fixa} $n$ i una quantitat $m$ que se suma (o es resta) per cada unitat: $y=mx+n$.}

\ex{1} Una espelma fa 24\,cm i es crema 3\,cm cada hora. Escriu la fórmula de l'altura $y$ al cap de $x$ hores i dibuixa'n el gràfic. Quant farà al cap de 5\,h? Quan s'acabarà?

\ex{1} Una tarifa de telèfon costa 6\,€ al mes més 0,08\,€ per minut. Quant paguem si parlem 50 minuts? Quants minuts hem parlat si paguem 9\,€?

\ex{2} Un dipòsit té 100\,L d'aigua i una aixeta hi aboca 25\,L cada minut. Escriu la fórmula. Quants litres hi haurà al cap de 6 minuts? Si hi caben 600\,L, quan s'omplirà?

\ex{2} Un lampista cobra segons aquesta taula. Troba la fórmula i calcula quant cobrarà per 6 hores. Quantes hores ha treballat si cobra 300\,€?\par\smallskip\begin{center}$\begin{array}{|c|ccc|}\hline \text{Hores}&1&2&3\\\hline \text{Preu (€)}&55&90&125\\\hline\end{array}$\end{center}

\ex{2} Per a passar de graus Celsius ($x$) a graus Fahrenheit ($y$) es fa servir la fórmula $y=1{,}8x+32$. A quants $^\circ$F equivalen 0\,$^\circ$C, 25\,$^\circ$C i 100\,$^\circ$C? A quants $^\circ$C equivalen 50\,$^\circ$F?

\ex{2} Un aparcament cobra 3\,€ per entrar i 2\,€ per cada hora. Escriu la fórmula i fes una taula per a 1, 2, 3, 4 i 5 hores. Quantes hores hi hem deixat el cotxe si paguem 19\,€?

\ex{3} Per a llogar una furgoneta, l'empresa A cobra 30\,€ més 0,25\,€ per km i l'empresa B, 50\,€ més 0,15\,€ per km. Quina convé si fem 150\,km? Per a quants km costen el mateix? Fes el gràfic de les dues.

\ex{3} El gimnàs A cobra 45\,€ de matrícula i 20\,€ al mes; el gimnàs B no té matrícula i cobra 35\,€ al mes. Fes el gràfic de les dues tarifes. Quan costen el mateix? Quin convé si vols anar-hi 6 mesos?

\vfill
\begin{center}\color{muted}\small\faCheckCircle[regular]\ \ Ànim! Si et costa algun exercici, repassa els apunts interactius de la unitat.\end{center}

\end{document}
